\chapter{Conception}
\label{chap:conception}

\begin{introchapitre}
Ce chapitre présente les choix de conception de SentinelOps. Il décrit l'architecture
globale, la conception de l'agent et de son automate d'alertes, la méthode de détection
d'anomalies, le modèle de données et les principales interactions sous forme de diagrammes
de séquence.
\end{introchapitre}

\section{Architecture globale}

SentinelOps s'exécute sur un seul hôte, sous forme de cinq conteneurs Docker. Deux font le
travail principal, deux protègent l'accès à Docker et un prépare le volume de données
(figure~\ref{fig:architecture}).

\begin{figure}[htbp]
\centering
\resizebox{\linewidth}{!}{% Architecture globale de SentinelOps (TikZ)
\begin{tikzpicture}[
  font=\sffamily\footnotesize,
  box/.style={draw=grisTexte, rounded corners=3pt, align=left, inner sep=5pt, text width=3.6cm, minimum height=1.55cm},
  main/.style={box, draw=bleuSentinel, very thick, fill=bleuSentinel!10},
  arr/.style={-{Stealth[length=2.2mm]}, draw=grisTexte, thick},
  lbl/.style={font=\sffamily\scriptsize, text=grisTexte, fill=white, inner sep=1pt}]
  % colonne 1 : sources et utilisateur
  \node[box] (host) at (0,0) {\textbf{Hôte}\\/proc et / de la machine\\montés en lecture seule};
  \node[box] (prom) at (0,-2.4) {\textbf{Exportateurs Prometheus}\\tout point /metrics\\(optionnel)};
  \node[box] (nav)  at (0,-4.8) {\textbf{Navigateur}\\interface React\\mise à jour toutes les 2 s};
  % colonne 2 : applications et base
  \node[box] (agent) at (5.7,0) {\textbf{agent} · Python 3.11\\collecte toutes les 5 s\\règles + anomalies\\automate d'états};
  \node[main] (db)   at (5.7,-2.4) {\textbf{SQLite (WAL)}\\métriques · alertes\\outbox · utilisateurs};
  \node[box] (dash)  at (5.7,-4.8) {\textbf{dashboard} · Node.js 24\\API REST + WebSocket\\authentification\\worker de notifications};
  \node[box, minimum height=0.9cm] (notif) at (5.7,-7.0) {\textbf{Slack · e-mail (SMTP)}};
  % colonne 3 : accès Docker filtré
  \node[box] (pro)  at (11.4,0) {\textbf{docker-proxy-ro}\\GET conteneurs seulement\\réseau interne};
  \node[box] (dock) at (11.4,-2.4) {\textbf{Moteur Docker}\\docker.sock\\= accès root};
  \node[box] (pops) at (11.4,-4.8) {\textbf{docker-proxy-ops}\\journaux + redémarrage\\réseau interne};
  % flèches
  \draw[arr] (host) -- node[lbl,above]{psutil} (agent);
  \draw[arr] (prom.east) -- ++(0.85,0) |- node[lbl,pos=0.25,right]{scrape} ([yshift=-4mm]agent.west);
  \draw[arr] (agent) -- node[lbl,right]{écrit} (db);
  \draw[arr,{Stealth[length=2.2mm]}-{Stealth[length=2.2mm]}] (dash) -- node[lbl,right]{lit, acquitte} (db);
  \draw[arr,{Stealth[length=2.2mm]}-{Stealth[length=2.2mm]}] (nav) -- node[lbl,above]{HTTP/WS} (dash);
  \draw[arr] (dash) -- node[lbl,right]{notifie} (notif);
  \draw[arr] (agent) -- node[lbl,above]{stats} (pro);
  \draw[arr] (dash) -- node[lbl,above]{restart} (pops);
  \draw[arr] (pro) -- (dock);
  \draw[arr] (pops) -- (dock);
  % cadre de l'hôte Docker
  \begin{scope}[on background layer]
    \node[draw=bleuSentinel, dashed, rounded corners=6pt, fit=(agent)(db)(dash)(pro)(dock)(pops), inner sep=9pt,
          label={[font=\sffamily\scriptsize\bfseries, text=bleuSentinel, anchor=south east]north east:Hôte Docker · \texttt{sen run}}] {};
  \end{scope}
\end{tikzpicture}
}
\caption{Architecture globale de SentinelOps}
\label{fig:architecture}
\end{figure}

\begin{table}[htbp]
\centering
\small
\begin{tabularx}{\linewidth}{L{3cm}L{3.2cm}Y}
\toprule
\textbf{Conteneur} & \textbf{Technologie} & \textbf{Responsabilité} \\
\midrule
\code{agent} & Python 3.11 (asyncio) & collecter, détecter, faire évoluer les alertes, nettoyer les anciennes données \\
\code{dashboard} & Node.js 24, Express, React & interface web, API REST, authentification, WebSocket, envoi des notifications \\
\code{docker-proxy-ro} & proxy HAProxy & n'autorise que la lecture des conteneurs pour l'agent \\
\code{docker-proxy-ops} & proxy HAProxy & n'autorise que les journaux et le redémarrage pour le dashboard \\
\code{volume-init} & shell (exécution unique) & attribue le volume à l'utilisateur non privilégié (uid 1000) \\
\bottomrule
\end{tabularx}
\caption{Rôle des conteneurs}
\label{tab:conteneurs}
\end{table}

\begin{aretenir}[Principe d'architecture]
L'agent et le dashboard ne s'appellent jamais directement : la \textbf{base SQLite est le
contrat} entre eux. L'agent écrit les métriques et les alertes, le dashboard les lit et
enregistre les actions des opérateurs. Chacun peut redémarrer sans interrompre l'autre.
\end{aretenir}

Les règles de conception suivantes ont guidé la refonte de la version 0.2 :

\begin{table}[htbp]
\centering
\small
\begin{tabularx}{\linewidth}{L{5cm}Y}
\toprule
\textbf{Règle} & \textbf{Problème résolu} \\
\midrule
Un seul schéma, défini par des migrations SQL versionnées partagées & le schéma défini à trois endroits en v0.1 : une colonne utilisée n'existait nulle part \\
Un seul écrivain par groupe de colonnes & concurrence et responsabilités floues \\
Outbox transactionnelle pour les notifications & notifications perdues ou envoyées pour une alerte annulée \\
Automate d'états avec durée et index unique & alertes bruyantes et dupliquées \\
Interface \code{Detector} commune & détecteur d'anomalies inutilisé ; extensibilité vers l'apprentissage automatique \\
API fermée par défaut & API entièrement ouverte en v0.1 \\
Accès Docker par proxies au moindre privilège & socket Docker monté en lecture-écriture \\
Horodatages en millisecondes UTC & formats de date mélangés \\
\bottomrule
\end{tabularx}
\caption{Règles de conception de la version 0.2}
\label{tab:regles-conception}
\end{table}

\section{Conception de l'agent}

\subsection{Chaîne de traitement}

L'agent exécute une boucle à cadence fixe de 5 secondes. Les opérations lentes (appels à
Docker, lecture des processus) s'exécutent dans des fils séparés afin de ne jamais bloquer la
boucle (figure~\ref{fig:pipeline}).

\diagramme[0.72\linewidth]{pipeline-agent.pdf}{Chaîne de traitement de l'agent}{fig:pipeline}

\subsection{Diagramme de classes}

Le contrat central est la classe \code{Signal} : un détecteur, quel qu'il soit, ne produit que
des signaux, et le gestionnaire d'alertes ne consomme que des signaux. Ajouter un nouveau
détecteur (par exemple un modèle Isolation Forest prévu en v0.2) ne demande donc aucune
modification du cycle de vie des alertes (figure~\ref{fig:classes}).

\diagramme[0.88\linewidth]{classes-agent.pdf}{Diagramme de classes de l'agent}{fig:classes}

\subsection{Automate d'états d'une alerte}

Le cycle de vie suit le modèle des règles d'alerte de Prometheus \cite{prometheus-rules}
(figure~\ref{fig:etats}).

\begin{figure}[htbp]
\centering
\resizebox{\linewidth}{!}{% Automate d'états d'une alerte (TikZ)
\begin{tikzpicture}[
  font=\sffamily\footnotesize,
  st/.style={draw=grisTexte, thick, rounded corners=4pt, minimum width=3.1cm, minimum height=1.25cm, align=center},
  fire/.style={st, draw=bleuSentinel, very thick, fill=bleuSentinel!10},
  arr/.style={-{Stealth[length=2.2mm]}, draw=grisTexte, thick},
  lbl/.style={font=\sffamily\scriptsize, text=grisTexte, align=center}]
  \fill[grisTexte] (-1.6,0) circle (1.6mm);
  \node[lbl, below=2mm] at (-1.6,0) {condition\\vraie};
  \node[st]   (p) at (1.6,0)  {\textbf{PENDING}\\[-1pt]{\scriptsize en mémoire seulement}};
  \node[fire] (f) at (7.6,0)  {\textbf{FIRING}\\[-1pt]{\scriptsize alerte + outbox}};
  \node[st]   (r) at (13.6,0) {\textbf{RESOLVED}\\[-1pt]{\scriptsize outbox « resolved »}};
  \node[draw=grisTexte, dashed, rounded corners=4pt, minimum width=2.8cm, minimum height=0.8cm, font=\sffamily\scriptsize] (d) at (-0.4,-2.3) {abandon silencieux};
  \draw[arr] (-1.4,0) -- (p);
  \draw[arr] (p) -- node[lbl, above=2pt]{maintenue pendant\\la durée « for »} (f);
  \draw[arr] (f) -- node[lbl, above=2pt]{disparue, règle supprimée\\ou résolution manuelle} (r);
  \draw[arr] ([xshift=-6mm]p.south) |- node[lbl, pos=0.3, left]{disparue\\avant « for »} (d.east);
  \draw[arr] (f.north) .. controls +(0.9,1.5) and +(-0.9,1.5) .. node[lbl, above]{toujours vraie : last\_seen, hit\_count + 1} (f.north east);
  \draw[arr] (r.south) -- ++(0,-2.6) -| node[lbl, pos=0.3, below]{résolue à la main alors que la condition est toujours vraie} ([xshift=8mm]p.south);
  \node[lbl, text width=4.6cm, below=0.5cm of f, align=center, text=bleuSentinel] {Acquitter = attribut \texttt{acked\_by}\\l'alerte reste FIRING et\\se résout toute seule};
\end{tikzpicture}
}
\caption{Automate d'états d'une alerte}
\label{fig:etats}
\end{figure}

\begin{itemize}
  \item \textbf{PENDING} : la condition vient d'apparaître. L'état est conservé
        \emph{uniquement en mémoire} : rien n'est écrit en base et personne n'est notifié. Si
        la condition disparaît avant la durée \code{for}, elle est abandonnée silencieusement :
        c'est ce qui supprime le bruit des pics passagers.
  \item \textbf{FIRING} : la condition a duré au moins \code{for} secondes. Une transaction
        unique écrit l'alerte et une ligne de notification. Les cycles suivants mettent
        seulement à jour \code{last_seen}, le compteur et les dernières preuves.
  \item \textbf{RESOLVED} : la condition a disparu (ou la règle a été supprimée). Une
        transaction marque l'alerte résolue (\code{resolved_by = auto}) et crée une
        notification de résolution.
  \item L'\textbf{acquittement} n'est pas un état mais un attribut (\code{acked_by},
        \code{acked_at}) : une alerte acquittée reste FIRING et se résout automatiquement.
\end{itemize}

Un \textbf{index unique partiel} sur l'empreinte (\code{sha256(règle:entité)}) des alertes à
l'état FIRING garantit au niveau de la base qu'une même règle ne peut pas avoir deux alertes
ouvertes sur une même entité, même en cas de concurrence.

\subsection{Corrélation en incidents}

Une nouvelle alerte rejoint l'incident d'une autre alerte ouverte si elles concernent la
\textbf{même entité} et ont commencé à moins de \textbf{120 secondes} d'écart ; sinon un nouvel
incident est créé avec un identifiant de la forme \code{inc-} suivi de 8 caractères
hexadécimaux. Par exemple, lors des tests, une alerte de seuil et une alerte d'anomalie sur le
CPU du même conteneur, apparues à la même seconde, ont été regroupées dans un seul incident.

\section{Détection d'anomalies}

\subsection{Principe : le score z robuste}

Le détecteur signale une valeur éloignée du comportement récent d'une métrique sur une entité,
sans seuil défini à la main. Pour chaque couple (entité, métrique), il conserve les 30 dernières
valeurs (2 min 30 s) et, à partir de 10 valeurs, évalue chaque nouvelle mesure $x$ :

\begin{equation}
  \tilde{x} = \operatorname{médiane}(\text{fenêtre}), \qquad
  \sigma = \max\!\left(1{,}4826 \times \operatorname{médiane}\left(\lvert x_i - \tilde{x}\rvert\right),\; 2\right), \qquad
  z = \frac{x - \tilde{x}}{\sigma}
\end{equation}

La valeur est anormale si $\lvert z \rvert > 3$ \textbf{et} si elle s'écarte d'au moins 8 unités
de la médiane. Le facteur 1,4826 rend l'écart absolu médian (MAD) comparable à un écart-type
pour des données normales \cite{leys2013,iglewicz1993}. Les deux garde-fous (plancher de
$\sigma$ et écart minimal) évitent les fausses alertes sur des métriques presque plates. La
plage attendue $\tilde{x} \pm 3\sigma$ est enregistrée avec chaque mesure et affichée sur les
courbes.

\subsection{Exemple de calcul}

Dix dernières valeurs du CPU d'un conteneur : 20, 21, 19, 20, 22, 20, 21, 19, 20, 21 ; nouvelle
valeur : 75.
\begin{enumerate}
  \item Médiane $\tilde{x} = 20$.
  \item Écarts à la médiane : 0, 1, 1, 0, 2, 0, 1, 1, 0, 1 ; leur médiane vaut 1, donc
        $1{,}4826 \times 1 = 1{,}48$, relevé au plancher $\sigma = 2$.
  \item $z = (75 - 20) / 2 = 27{,}5 > 3$ et l'écart de 55 dépasse 8 : la valeur est anormale.
  \item Une alerte n'est créée que si l'anomalie dure 15 secondes (3 cycles), exactement comme
        la durée \code{for} d'une règle.
\end{enumerate}

\subsection{Pourquoi la médiane plutôt que la moyenne}

Avec la moyenne et l'écart-type, une anomalie durable \og se cache \fg{} : le saut lui-même
augmente l'écart-type, ce qui fait chuter le score dès les mesures suivantes. La médiane et le
MAD ne sont presque pas influencés par quelques valeurs extrêmes. Ce choix a été validé par une
mesure présentée au chapitre~\ref{chap:devsecops} (section~\ref{sec:evaluation}).

\section{Modèle de données}

\diagramme[0.92\linewidth]{modele-donnees.pdf}{Modèle de données (principales tables)}{fig:mcd}

Toutes les dates sont des entiers en millisecondes UTC. La base fonctionne en mode WAL
(\emph{write-ahead logging}), qui permet au dashboard de lire pendant que l'agent écrit
\cite{sqlite-wal}.

\begin{table}[htbp]
\centering
\small
\begin{tabularx}{\linewidth}{L{3.6cm}L{4.2cm}Y}
\toprule
\textbf{Table} & \textbf{Écrite par} & \textbf{Conservation} \\
\midrule
\code{metrics} & agent & 24 heures \\
\code{metrics_1m} & agent (agrégats par minute) & 7 jours \\
\code{alerts} & agent (détection) et dashboard (acquittement, résolution manuelle) & 30 jours après résolution \\
\code{notifications_outbox} & agent et dashboard, dans la même transaction que l'alerte & 7 jours après traitement \\
\code{rules}, \code{settings} & dashboard (initialisation des règles par l'agent) & permanente \\
\code{users}, \code{sessions} & dashboard & sessions expirées supprimées \\
\code{schema_version}, \code{app_meta} & les deux / agent & permanente \\
\bottomrule
\end{tabularx}
\caption{Propriété et conservation des tables}
\label{tab:tables}
\end{table}

\subsection{Migrations versionnées}

Le schéma est défini une seule fois, dans des fichiers SQL numérotés
(\nolinkurl{shared/migrations/001_init.sql}). L'agent (Python) et le dashboard (Node.js) exécutent
les mêmes fichiers au démarrage, sous un verrou d'écriture exclusif (\code{BEGIN IMMEDIATE}) :
celui qui démarre en premier applique les migrations manquantes, l'autre constate qu'il n'y a
rien à faire. La migration 001 met également à niveau une base de la version 0.1 en conservant
les utilisateurs, les règles et les paramètres.

\subsection{Outbox transactionnelle}

Enregistrer une alerte et envoyer sa notification sont deux actions distinctes : un arrêt entre
les deux ferait perdre la notification. Le patron \emph{transactional outbox}
\cite{outbox} résout ce problème : l'alerte et une ligne de notification sont écrites dans la
même transaction, puis un processus indépendant livre les notifications en attente et réessaie
en cas d'échec. La livraison est dite \og au moins une fois \fg{} : un doublon est possible
après un échec partiel, une perte ne l'est pas.

\section{Diagramme de séquence du cycle de détection}

\diagramme[0.85\linewidth]{seq-detection.pdf}{Diagramme de séquence du cycle de détection}{fig:seq-detection}


L'authentification (vérification de l'empreinte, création d'une session, ouverture du WebSocket avec le jeton) et l'acquittement suivent le même principe : le dashboard écrit l'action de l'opérateur dans la base, que l'agent prend en compte au cycle suivant.

\begin{conclusionchapitre}
La conception repose sur un contrat simple entre deux programmes (la base de données), un
cycle de vie des alertes explicite et une détection extensible. Le chapitre suivant présente sa
réalisation.
\end{conclusionchapitre}
